\FloatBarrier
\section{Introduction}
\label{sec:introduction}
\FloatBarrier
\subsection{Motivation}
Algebraic tensor rewriting is used to optimize tensor computations by replacing
one computation graph with an equivalent graph. The equation describing a
rewrite can be short even when its implementation is complex. For example,
reassociating a product from $(XY)Z$ to $X(YZ)$ changes only the nesting of two
dots in an algebraic expression. For generalized tensor dots, however, an
implementation must also reconstruct which dimensions are batched or
contracted, preserve their pairings, and arrange the output dimensions in the
required order. XLA's \code{AssociativeReorderNestedDot} routine illustrates
these tasks: it constructs new dimension descriptors and, when needed, an
output permutation, alongside applicability and cost checks~\cite{xla-simplifier}.
This observation concerns the kinds of work required, rather than a measured
line-count comparison. \todo{Classify the pinned XLA routine's attribute logic
separately from legality, profitability, and graph construction before
reporting its size.}

Symbolic rewriting frameworks such as egg support concise operation
patterns~\cite{egg}. They can represent operation attributes, but their pattern
matching does not itself derive attribute transformations from tensor
semantics. For a tensor rewrite, the user must supply the additional
expressions or custom host-language logic; egg's \code{Applier} interface, for
example, permits application code to inspect analysis data and construct a
replacement~\cite{egg-applier}. Merely matching the operations and attributes
on the left-hand side (LHS) does not determine the attributes needed on the
right-hand side (RHS).

Manual attribute transformation is a major source of this complexity, although
legality checks, numerical policy, and profitability also contribute. Cascaded
operations change axis positions, operand pairings, and output order. Equal
dimension sizes do not establish that two dimensions denote the same logical
coordinate. Moreover, the transformation should work for tensors of arbitrary
rank; working out a few concrete matrix or fixed-rank examples does not
establish the general relationship.

Our research question is therefore: \emph{What if we could design a language
that makes tensor rewriting as easy and intuitive as writing its algebraic
form, while leaving complex attribute transformation to the compiler?}
TEPL pursues this goal through operation semantics, symbolic axis mapping,
and automatic RHS attribute inference under explicit conditions.
Figure~\ref{fig:why-tepl} compares the authoring choices for the same
reassociation. A bare algebraic pattern supplies the replacement graph but
leaves both replacement descriptors unspecified. Custom implementation code
supplies them through handwritten dimension logic. TEPL supports both
user-written descriptor expressions and compiler inference from reusable
operation semantics. The example is unbatched and presupposes arithmetic
that permits reassociation; it does not establish that every generalized-dot
reassociation is valid.

\paperfigure{1}{figures/Why TEPL_ From Algebraic Rules to Practical Tensor Rewrites.pdf}{TEPL combines concise tensor rewrite patterns with explicit or inferred replacement attributes.}{fig:why-tepl}{Four panels compare an algebraic pattern, custom implementation logic,
explicit TEPL attribute expressions, and inferred TEPL attributes for the
same nested-dot replacement graph.}

\FloatBarrier
\subsection{Central Thesis}
TEPL describes tensor rewrites, their conditions, and their attribute changes
together. Users can write attribute changes directly or ask the compiler to
infer them from operation semantics. The compiler checks supported semantic
rewrites for validity and equivalent results.

The language lets users express algebraic patterns, applicability conditions,
and attribute changes in one place. It supports user-defined dialects, each
with its own operations, attributes, and semantics. An operation is defined
once within a dialect and then imported and reused across rules. Tensor
operations, graph patterns, attribute bindings, metadata conditions, and
replacement-attribute expressions are represented directly in the language.

The central technical advance pursued here is inference of omitted RHS
attributes. Users may also request this inference explicitly with
\code{auto}. Symbolic axis mapping constructs and checks reusable attribute
programs for arbitrary-rank tensors within the supported semantics. The
compiler reasons from operation legality and the user's stated conditions;
when those assumptions do not establish correctness, the semantic rule is
rejected as unproved. Inference does not silently add conditions to make a
proposed rewrite succeed.

\FloatBarrier
\subsection{Proposed Contributions}
This paper develops three connected contributions:
\begin{enumerate}
\item \textbf{A language for tensor rewrites.} TEPL supports user-defined
dialects and operations and expresses rewrite patterns, applicability
conditions, and attribute changes together. Users can capture and reuse
descriptors or compute replacement attributes directly in the language.
\item \textbf{Automatic attribute inference.} The compiler synthesizes omitted
RHS attributes for supported rewrites as reusable functions of LHS metadata,
without fixing tensor rank. Its Axes IR builds on prior aggregated-axis
representations and adds explicit ordered interfaces, operand mappings,
and value semantics for this synthesis task.
\item \textbf{Symbolic axis mapping for composed operations.} The method tracks
logical-axis correspondence across composed operations, splits and composes
arbitrary-length groups while preserving coordinate order and pairings,
and reuses established producer--consumer mappings for algebraic matching
and attribute construction.
\end{enumerate}
These are the design and mechanism claims of this draft. The correctness
contract is stated as an intended guarantee, and the practical benefits are
questions for the planned evaluation, rather than established experimental
results.
